\subsection{Malliavin representation}\label{subsec: malliavin3F}
Following \cite{guyonLabordere}, we give an alternative expression to the calibration condition (\ref{eq: threeFactorCalibrationCondition}), in order to achieve a more accurate and smooth estimation of the local volatility for strikes that are far from the money.

\subsubsection{The general formula}
Two elementary facts about the forward measures are needed.
\begin{lem}[Forward measures]\label{lem: forwardMeasures}
For every $t$, with $F_{0t}=S_0P^f_{0t}/P^d_{0t}$ the forward FX rate,
\begin{equation}\label{eq: RNforward}
\left.\frac{d\mathbb{Q}^{f^t}}{d\mathbb{Q}^{d^t}}\right|_{\mathcal{F}_t}=\frac{S_t}{F_{0t}} .
\end{equation}
Consequently, if $q^{d^t}$ and $q^{f^t}$ denote the densities of $S_t$ under the two measures,
\begin{itemize}
\item[(i)] $q^{f^t}(K)=\frac{K}{F_{0t}}\,q^{d^t}(K)$, and therefore $S_0P^f_{0t}\,q^{f^t}(K)=K\,P^d_{0t}\,q^{d^t}(K)=K\,\partial_K^2C(t,K)$;
\item[(ii)] $\mathbb{E}^{f^t}[H\,|\,S_t=K]=\mathbb{E}^{d^t}[H\,|\,S_t=K]$ for every integrable $\mathcal{F}_t$-measurable $H$.
\end{itemize}
\end{lem}
\begin{proof}
In domestic currency the numeraire of $\mathbb{Q}^{f^t}$ is $S_uP^f_{ut}$ and the numeraire of $\mathbb{Q}^{d^t}$ is $P^d_{ut}$; their ratio at $u=t$ is $S_t$, and $F_{0t}$ is its value at $u=0$. Since the density (\ref{eq: RNforward}) is a function of $S_t$ only, (i) is immediate and (ii) follows from Bayes' formula.
\end{proof}

\paragraph{Convention.} Write $(W^S,B^d,B^f)^\top=\mathcal{C}Z$ with $Z$ a three-dimensional standard Brownian motion and $\mathcal{C}\mathcal{C}^\top=R$. The Malliavin derivative with respect to the rate driver $B^d$ is $\MD^{B^d}_sF:=\langle \MD^Z_sF,\mathcal{C}^\top e_d\rangle$: it is the sensitivity of $F$ to a perturbation of the path of $B^d$ which moves $W^S$ and $B^f$ by the regression coefficients $\rho_{Sd}$ and $\rho_{df}$. With this convention $\mathbb{E}[F\int_0^th(s)dB^d_s]=\mathbb{E}[\int_0^th(s)\MD^{B^d}_sF\,ds]$. The same convention is used for the drivers of the other measures.

For this single-driver representation assume each rate is autonomous under its own forward measure and measurable with respect to its own rate-driver filtration; otherwise a full vector Clark--Ocone representation is required. By the Clark-Ocone formula, $r^d_t-r^0_d=\int_0^t\sigma^{d^t}_r(s)\,dB^{d^t}_s$ under $\mathbb{Q}^{d^t}$ and $r^f_t-r^0_f=\int_0^t\sigma^{f^t}_r(s)\,d\widetilde{B}^{f^t}_s$ under $\mathbb{Q}^{f^t}$, where $\sigma^{d^t}_r(s)=\mathbb{E}^{d^t}[\MD^{B^{d^t}}_sr^d_t|\mathcal{F}_s]$ and $\sigma^{f^t}_r(s)=\mathbb{E}^{f^t}[\MD^{\widetilde{B}^{f^t}}_sr^f_t|\mathcal{F}_s]$. The derivation of section 11.7.6.1 in \cite{guyonLabordere}, applied to each term of (\ref{eq: threeFactorCalibrationCondition}) under its own forward measure and combined with Lemma \ref{lem: forwardMeasures} (i), gives
\begin{align*}
\sigma(t,K)^2 = \sigma_{Dup}(t,K)^2 & - \frac{2}{K}\int _0^t \mathbb{E}^{d^t}\big[\sigma_r^{d^t}(s)\, \MD _s^{B^{d^t}}S_t  \big|S_t = K \big]ds \\
& +  \frac{2}{K}\int _0^t \mathbb{E}^{f^t}\big[\sigma_r^{f^t}(s)\, \MD _s^{\widetilde{B}^{f^t}}S_t  \big|S_t = K \big]ds
\numberthis \label{eq: threeFactorGeneralMalliavin}
\end{align*}
Both terms carry the factor $2/K$. In the foreign term,
$$S_0\,\mathbb{E}^f\big[D^f_{0t}\,\delta_K(S_t)\,\cdot\,\big]=S_0P^f_{0t}\,q^{f^t}(K)\,\mathbb{E}^{f^t}\big[\,\cdot\,\big|S_t=K\big],$$
and the density of $S_t$ under the foreign forward measure is $K\partial_K^2C/(S_0P^f_{0t})$, not $\partial^2_KC/P^f_{0t}$.

For Hull-White rates the Clark-Ocone integrands are deterministic,
\begin{equation}\label{eq: kernels}
\sigma_r^{d^t}(s)=\sigma^d_r\,e^{-\kappa_d(t-s)},\qquad \sigma_r^{f^t}(s)=\sigma^f_r\,e^{-\kappa_f(t-s)},
\end{equation}
because each rate is autonomous and Gaussian under its own forward measure. What remains is to evaluate the two integrated sensitivities of $S_t$.

\subsubsection{Malliavin derivative demystified}
Before moving to the actual computation of the Malliavin derivatives in (\ref{eq: threeFactorGeneralMalliavin}) required for an explicit formulation of the Particle Algorithm in the case of three-factor FX SIR-LV model, I attempt to give a short introduction to the tools needed to understand the following derivations.
I need to assume familiarity with the basics of Malliavin calculus. These can be found in the first chapter of \cite{nualart} and I will refer to the Nualart's book through out the computations.

The results I am going to use are the \textit{commutative relation} between the Derivative and Divergence operators (Proposition 1.3.8 in \cite{nualart}) and the \textit{chain rule} for the Derivative operator (Proposition 1.2.3 in \cite{nualart}).

\begin{prop}[Commutative Relation between $\MD$ and $\delta$]\ \\
Provided that the condition of Proposition 1.3.8 in \cite{nualart} are satisfied, one has:

\begin{equation}\label{eq: commumativeRelation}
\MD_s( \delta(u) )= u_s\mathbb{I}_{s\leq t} + \int_0^t\MD_s u_r \mathbb{I}_{s\leq t}dW_r
\end{equation}
for a stochastic process $u$ on $L^2([0,t] \times \Omega)$.
\end{prop}

To get familiar with such tools, I am going to solve the following easy exercise (exercise 2.2.1 in\cite{nualart} ):\\

\begin{exe}
Let $\sigma$ and $b$ be continuously differentiable functions on $\mathbb{R}$ with bounded
derivatives. Consider the solution $X = \{X_t, t \in [0, T]\}$ of the stochastic
differential equation
$$ X_t = x_0 + \int_0^t\sigma(X_r)dW_r + \int_0^tb(X_r)dr$$
Show that for $s\leq t$ we have
$$ \MD_sX_t = \sigma(X_s)\exp\bigg(\int_s^t\sigma'(X_r)dW_r  + \int_s^t[b'(X_r)-\frac{1}{2}\sigma'(X_r)^2]dr  \bigg) $$
Where $\MD$ is the Malliavin derivative with respect to the Brownian Motion $W$.
\end{exe}\
\\
\begin{sol}
Let's first rewrite the previous SDE as:
$$ X_t = x_0 + \delta\big(\sigma(X)\big) + <b(X),\mathbb{I}_{[0,t]}> $$
where $\delta$ is the Skorohod integral and $<\cdot, \cdot>$ is the inner product, both under the White Noise case (Section 1.2.1 in \cite{nualart}).\\
Then:
\begin{align*}
\MD_sX_t = & \MD_s\delta\big(\sigma(X)\big) + <\MD_sb(X),\mathbb{I}_{[0,t]}>\\
= & \sigma(X_s)\mathbb{I}_{s\leq t} + \delta\big(\MD_s\sigma(X)\mathbb{I}_{s\leq t}\big) +  <\MD_sb(X),\mathbb{I}_{[0,t]}>
\end{align*} 
Where we get the second line from an application of the \textit{commutative relation}.
By applying the \textit{chain rule} of the Malliavin derivative to $\sigma$ and $b$ we get:
\begin{align*}
\MD_sX_t =\MD_sX_t\mathbb{I}_{s\leq t}= & \sigma(X_s)\mathbb{I}_{s\leq t} + \delta\big( \sigma'(X)\cdot \MD_sX \mathbb{I}_{s\leq t}\big) + < b'(X) \cdot\MD_s X,\mathbb{I}_{[0,t]}>\\
= & \sigma(X_s) + \int_0^t \sigma'(X_r)\cdot\MD_sX_r \mathbb{I}_{s\leq t} dW_r + \int_0^t b'(X_r) \cdot\MD_s X_r\mathbb{I}_{s\leq t} dr
\end{align*} 
Define now $S_t:= \MD_sX_t\mathbb{I}_{s\leq t}$, one may write
$$ S_t = \sigma(X_s)\mathbb{I}_{s\leq t} + \int_0^t \sigma'(X_r)S_rdW_r + \int_0^t b'(X_r) S_rdr,$$
thus
$$ dS_t = \sigma'(X_t)S_tdW_t + b'(X_t) S_t dt$$
which is the SDE driving a Geometric Brownian Motion with initial value equal to $\sigma(X_s)$. So the result is proven $\diamond$ .
\end{sol}

The exercise is the scalar case of a general fact used below: for a diffusion $X$, $\MD_sX_t=J_tJ_s^{-1}\times$(diffusion coefficient at $s$), where $J$ is the first variation (tangent) process of $X$. For a system, $J$ is matrix-valued.

\subsubsection{Why the foreign rate cannot be treated by a change of measure}\label{subsubsec: obstacle}
Under $\mathbb{Q}^{f}$ the foreign rate is autonomous, exactly as the domestic rate is under $\mathbb{Q}^d$, and it is tempting to compute its Malliavin derivative there and to transport the result to the domestic measure (a change-of-variable formula for the Malliavin derivative under a Girsanov shift is given in \cite{ewald}). This does not lead to a closed form, for two reasons.

First, the domestic term of (\ref{eq: threeFactorGeneralMalliavin}) requires the derivative with respect to $B^{d^t}$, not with respect to the foreign driver.
Second, and more fundamentally, by (\ref{measureChange: domesticToForeign}) the drivers of the two measures differ by $\rho_{Sf}\int_0^t\sigma(u,S_u)du$, which is a functional of the path: a perturbation of the noise moves $S$, hence the Girsanov kernel, hence $r^f$. One Brownian motion cannot be substituted for the other inside a Malliavin derivative. Indeed, differentiating the Hull-White foreign rate directly under $\mathbb{Q}^{d^T}$ gives
$$ d\,\MD^{B^{d^T}}_sr^f_t=\Big[-\kappa_f\,\MD^{B^{d^T}}_sr^f_t-\rho_{Sf}\,\sigma^f_r\,\partial_S\sigma(t,S_t)\,\MD^{B^{d^T}}_sS_t\Big]dt, $$
while $\MD^{B^{d^T}}_sS_t$ depends on $\MD^{B^{d^T}}_sr^f_t$ through the drift $r^d_t-r^f_t$ of the spot. The pair $(\MD_sS_t,\MD_sr^f_t)$ solves a coupled linear system with random coefficients; it has no scalar factorisation of the type $V_t/V_s$ that made the one-factor computation of Section \ref{subsec: malliavin1F} explicit. The symmetric statement holds under the foreign measure, where it is the domestic rate that carries a quanto drift $+\rho_{Sd}\sigma^d_r\sigma(t,S_t)$.

The way out is not to solve for the two-parameter family $\MD_sS_t$, $0\leq s\leq t$, at all: because the kernels (\ref{eq: kernels}) factorise as $e^{-\kappa t}\cdot e^{\kappa s}$, the $s$-integral in (\ref{eq: threeFactorGeneralMalliavin}) can be carried through the linear system.

\subsubsection{Integrated sensitivities}
Let $X_t=(S_t,r^d_t,r^f_t)^\top$. Under $\mathbb{Q}^d$, $dX_t=b(t,X_t)dt+\mathcal{A}(t,X_t)dZ_t$ with
$$ b(t,X)=\begin{pmatrix} S(r^d-r^f)\\ \kappa_d(\theta_d(t)-r^d)\\ \kappa_f(\theta_f(t)-r^f)-\rho_{Sf}\sigma^f_r\,\sigma(t,S)\end{pmatrix},\qquad
\mathcal{A}(t,X)=\begin{pmatrix} S\sigma(t,S)&0&0\\ 0&\sigma^d_r&0\\ 0&0&\sigma^f_r\end{pmatrix}\mathcal{C}. $$
Let $M_t=\nabla_Xb(t,X_t)$ and $G_{j,t}=\nabla_X\mathcal{A}_{\cdot j}(t,X_t)$ be the Jacobians of the drift and of the columns of $\mathcal{A}$, and define the integrated sensitivities
\begin{equation}\label{eq: integratedSensitivities}
Y^d_t:=\int_0^t\sigma^d_re^{-\kappa_d(t-s)}\,\MD^{B^{d^t}}_sX_t\,ds,\qquad
Y^f_t:=\int_0^t\sigma^f_re^{-\kappa_f(t-s)}\,\MD^{\widetilde{B}^{f^t}}_sX_t\,ds .
\end{equation}
\paragraph{Standing hypotheses for the Hull--White representation.}\label{hyp: representation}
Fix a finite horizon $T_*>0$, deterministic initial states as in
(\ref{eq: centredStates}), and a positive semidefinite correlation matrix $R$.
Choose any square root $\mathcal C\mathcal C^\top=R$; its inverse is not used.
The rate parameters are constant with $\kappa_i>0$. Assume:
\begin{itemize}
\item[(H1)] The coefficient $\sigma(t,S)$ is deterministic, continuous in $t$
and locally continuously differentiable in $S>0$. The state has a nonexplosive
strong solution with $S_t>0$. Its Cameron--Martin first variations exist under
the domestic and foreign risk-neutral measures and solve the differentiated
state equations. No global bound on $S\partial_S\sigma$ is imposed.
\item[(H2)] All numeraire-change stochastic exponentials are true martingales
on $[0,T_*]$. For each $t>0$, $S_t$ belongs to $\mathbb D^{1,2}$ under both
own $t$-forward measures. Its Malliavin derivative is square integrable on
$[0,t]\times\Omega$ under those measures, so integration against the bounded
Hull--White kernels and Malliavin duality are legitimate. The integrated
first variations and the drift and diffusion terms of their equations are
locally integrable to the orders needed for stochastic Fubini. In particular,
$\mathbb E^{d^t}|S_t\ell^d_t|+
\mathbb E^{f^t}|S_t\ell^f_t|<\infty$.
\item[(H3)] For the classical calibration statement, the discounted Tanaka
martingale is a true martingale and the rate-tail expectations are finite.
The model call surface is $C^{1,2}$ for $t>0,K>0$, and
$m(t,K)=\partial_{KK}C_m(t,K)>0$. The locally finite signed weighted measures
$\mathbb E^d[D^d_{0t}\ell^i_t\mathbf1_{S_t\in dK}]$ have continuous densities
$n_i(t,K)$. Conditional expectations at a strike use the continuous versions
$n_i(t,K)/m(t,K)$. These assumptions justify the indicator limit below.
\end{itemize}
The first two hypotheses suffice for the integrated system and smooth-payoff
duality. Hypothesis (H3) is needed only for the pointwise classical calibration
formula. Local differentiability alone is not asserted to imply (H2) or (H3).

\begin{prop}[Integrated sensitivities]\label{prop: integratedSensitivities}
Under the standing hypotheses, the processes $Y^d$ and $Y^f$ solve, with $Y^d_0=Y^f_0=0$,
\begin{align*}
dY_t^d &=\big[(M_t-\kappa_d I)Y_t^d+a_d(t,X_t)\big]dt +\sum_{j=1}^3G_{j,t}Y_t^d\,dZ_t^j \numberthis\label{eq: matrixSystem}\\
dY_t^f &=\big[(M_t+H_t-\kappa_f I)Y_t^f+a_f(t,X_t)\big]dt +\sum_{j=1}^3G_{j,t}Y_t^f\,dZ_t^j
\end{align*}
where $Z$ is the $\mathbb{Q}^d$-Brownian motion and
$$ a_d=\sigma^d_r\begin{pmatrix}\rho_{Sd}S\sigma\\ \sigma^d_r\\ \rho_{df}\sigma^f_r\end{pmatrix},\qquad
a_f=\sigma^f_r\begin{pmatrix}\rho_{Sf}S\sigma\\ \rho_{df}\sigma^d_r\\ \sigma^f_r\end{pmatrix},\qquad
H_t=\partial_S\sigma(t,S_t)\begin{pmatrix}S_t\sigma(t,S_t)\\ \rho_{Sd}\sigma^d_r\\ \rho_{Sf}\sigma^f_r\end{pmatrix}\begin{pmatrix}1&0&0\end{pmatrix}. $$
\end{prop}
\begin{proof}
For a deterministic $h$, $\int_0^th(s)\MD^{B}_sX_t\,ds$ is the derivative at $\varepsilon=0$ of the solution $X^\varepsilon_t$ obtained by shifting the driver $B$ by $\varepsilon\int_0^\cdot h(s)ds$ (and the correlated drivers accordingly): the Malliavin derivative in a Cameron-Martin direction is a directional derivative on the path space. Differentiating the SDE in $\varepsilon$ gives a linear SDE, the first variation equation with a source term.

\textit{Domestic.} For Hull-White rates $\mathbb{Q}^d$, $\mathbb{Q}^{d^t}$ and $\mathbb{Q}^{d^T}$ differ by deterministic drifts, (\ref{measureChange: domesticForwardToDomestic}), so shifting $B^{d^t}$ is the same as shifting the $\mathbb{Q}^d$-driver: $Z\to Z+\varepsilon\int h\,\mathcal{C}^\top e_d\,ds$. With $h(s)=\sigma^d_re^{\kappa_ds}$, the derivative $\mathcal{Y}_t=\partial_\varepsilon X^\varepsilon_t$ solves $d\mathcal{Y}_t=[M_t\mathcal{Y}_t+e^{\kappa_dt}a_d]dt+\sum_jG_{j,t}\mathcal{Y}_tdZ^j_t$, since $\mathcal{A}\,\mathcal{C}^\top e_d\,\sigma^d_r=a_d$. As $Y^d_t=e^{-\kappa_dt}\mathcal{Y}_t$, the first equation follows.

\textit{Foreign.} Again $\mathbb{Q}^f$ and $\mathbb{Q}^{f^t}$ differ by a deterministic drift, so the shift acts on the $\mathbb{Q}^f$-driver $\widetilde{Z}$, where by (\ref{measureChange: domesticToForeign}) $dZ_t=d\widetilde{Z}_t+\mathcal{C}^\top e_S\,\sigma(t,S_t)dt$. Under $\mathbb{Q}^f$ the drift of $X$ is $b+\sigma\,w$ with $w:=\mathcal{A}\,\mathcal{C}^\top e_S=(S\sigma,\rho_{Sd}\sigma^d_r,\rho_{Sf}\sigma^f_r)^\top$, so that the first variation equation has drift Jacobian $M_t+\nabla_X(\sigma w)$ and noise $\sum_jG_{j,t}\mathcal{Y}_td\widetilde{Z}^j_t=\sum_jG_{j,t}\mathcal{Y}_tdZ^j_t-\sigma\,\nabla_Xw\,\mathcal{Y}_tdt$. The two contributions of the Girsanov kernel combine into
$$ \nabla_X(\sigma w)-\sigma\nabla_Xw=w\,(\nabla_X\sigma)^\top=H_t . $$
The source term is $e^{\kappa_ft}a_f$ and $Y^f_t=e^{-\kappa_ft}\mathcal{Y}_t$.
\end{proof}
The matrix $H_t$ is the diffusion direction of the spot driver times the gradient of the Girsanov kernel between $\mathbb{Q}^d$ and $\mathbb{Q}^f$. It is exactly what a substitution of Brownian motions loses, and it is the reason why $Y^f$ does not satisfy the same variational system as $Y^d$. The systems (\ref{eq: matrixSystem}) have no index $s$ left, and do not depend on a terminal maturity.

\subsubsection{Reduction to four scalar processes}
Two components of (\ref{eq: matrixSystem}) are deterministic: $Y^{d,d}_t=v_d(t)$ and $Y^{f,f}_t=v_f(t)$, where
$$ v_i(t)=(\sigma^i_r)^2\,\frac{1-e^{-2\kappa_it}}{2\kappa_i},\qquad i\in\{d,f\}, $$
is the variance of the Hull-White state $x^i_t$ (for $Y^{f,f}$, note that the third row of $M_t+H_t$ is $(0,0,-\kappa_f)$: the quanto term cancels, as it must under $\mathbb{Q}^f$). For the remaining four, pass to logarithmic sensitivities of the spot,
$$ \ell^d_t:=\frac{Y^{d,S}_t}{S_t},\qquad \zeta^f_t:=Y^{d,f}_t,\qquad \ell^f_t:=\frac{Y^{f,S}_t}{S_t},\qquad \zeta^d_t:=Y^{f,d}_t , $$
and let $g(t,S)=S\,\partial_S\sigma(t,S)$. By It\^o's formula the terms $r^d-r^f$ and $\sigma\,dW^S$ of the spot cancel and (\ref{eq: matrixSystem}) becomes
\begin{align*}
d\ell^d_t &= \ell^d_t\Big[g(t,S_t)\big(dW^S_t-\sigma(t,S_t)dt\big)-\kappa_d\,dt\Big]
 +\Big[v_d(t)-\zeta^f_t+\rho_{Sd}\,\sigma^d_r\,\sigma(t,S_t)\Big]dt \\
d\zeta^f_t &= \Big[-(\kappa_d+\kappa_f)\zeta^f_t-\rho_{Sf}\,\sigma^f_r\,g(t,S_t)\,\ell^d_t
 +\rho_{df}\,\sigma^d_r\sigma^f_r\Big]dt \numberthis \label{eq: sensitivities}\\
d\ell^f_t &= \ell^f_t\Big[g(t,S_t)\,dW^S_t-\kappa_f\,dt\Big]
 +\Big[\zeta^d_t-v_f(t)+\rho_{Sf}\,\sigma^f_r\,\sigma(t,S_t)\Big]dt \\
d\zeta^d_t &= \Big[-(\kappa_d+\kappa_f)\zeta^d_t+\rho_{Sd}\,\sigma^d_r\,g(t,S_t)\,\ell^f_t
 +\rho_{df}\,\sigma^d_r\sigma^f_r\Big]dt
\end{align*}
with $\ell^d_0=\zeta^f_0=\ell^f_0=\zeta^d_0=0$ and $W^S$ the $\mathbb{Q}^d$-Brownian motion; under $\mathbb{Q}^{d^T}$ substitute $dW^S_t=dW^{S^T}_t-\rho_{Sd}\sigma^d_rA^d(t,T)dt$.

Two features of (\ref{eq: sensitivities}) are worth noting. The domestic pair has $g\,(dW^S-\sigma dt)$ where the foreign pair has $g\,dW^S$: this is the first component of $H_t$. The cross terms $-\rho_{Sf}\sigma^f_r\,g\,\ell^d$ and $+\rho_{Sd}\sigma^d_r\,g\,\ell^f$ are the feedback of the spot sensitivity into the rate that carries the quanto drift: this is the coupling of Section \ref{subsubsec: obstacle}. When the local volatility has no skew ($g=0$) everything is deterministic and $\zeta^f_t=\zeta^d_t$ is the covariance of $x^d_t$ and $x^f_t$.

\begin{lem}[Smooth-payoff duality]\label{lem: duality}
Under (H1)--(H2), for each bounded $C^1$ function $\psi$ with bounded derivative,
\begin{align*}
\mathbb E^d[D^d_{0t}c^d_t\psi(S_t)]
 &=\mathbb E^d[D^d_{0t}S_t\ell^d_t\psi'(S_t)],\\
\mathbb E^d[D^d_{0t}S_tc^f_t\psi(S_t)]
 &=\mathbb E^d[D^d_{0t}S_t^2\ell^f_t\psi'(S_t)],
\qquad c^i_t=r^i_t-r_i^0(t).
\end{align*}
No strike density or strict positive definiteness of $R$ is required here.
\end{lem}
\begin{proof}
Under the domestic $t$-forward measure,
$c_t^d=\int_0^t\sigma_r^d e^{-\kappa_d(t-s)}dB_s^{d^t}$.
The chain rule and duality against this deterministic square-integrable kernel
show that $\mathbb E^{d^t}[c_t^d\psi(S_t)]
=\mathbb E^{d^t}[S_t\ell_t^d\psi'(S_t)]$.
Multiplication by $P^d_{0t}$ gives the first identity. Repeat the argument
under the foreign $t$-forward measure and multiply by $S_0P^f_{0t}$.
The identity $S_0\mathbb E^f[D^f_{0t}H]=\mathbb E^d[D^d_{0t}S_tH]$
then gives the second identity. All changes of measure are applied to the
already-defined path functional $\ell^f$, not to an untransformed Malliavin operator.
\end{proof}

\begin{prop}[Malliavin representation, Hull-White rates]\label{prop: main}
Let both rates follow Hull-White dynamics and impose (H1)--(H3).
Let the market call surface $C$ have positive strike density and a positive
Dupire variance, with $C(0,K)=(S_0-K)^+$, $C(t,0)=S_0P^f_{0t}$ and
$C(t,K)\to0$ as $K\to\infty$. Assume uniqueness for its linear forward Dupire
PDE in a solution class containing both $C$ and the model call surface $C_m$,
with these initial and boundary conditions. Then the existing three-factor
model is calibrated to $C$ if and only if
\begin{equation}\label{eq: mainFormula}
\boxed{\;\sigma(t,K)^2=\sigma_{Dup}(t,K)^2-2\,\mathbb{E}^{d^t}\big[\ell^d_t-\ell^f_t\,\big|\,S_t=K\big]\;}
\end{equation}
where $\ell^d,\ell^f$ solve (\ref{eq: sensitivities}). Equivalently, for any $T\geq t$,
\begin{align*}
\sigma(t,K)^2&=\sigma_{Dup}(t,K)^2
-2\,\frac{\mathbb{E}^{\mathbb{Q}^{d^T}}\big[(P^d_{tT})^{-1}(\ell^d_t-\ell^f_t)\,\big|\,S_t=K\big]}{\mathbb{E}^{\mathbb{Q}^{d^T}}\big[(P^d_{tT})^{-1}\,\big|\,S_t=K\big]} \numberthis\label{eq: mainFormulaForward}\\
&=\sigma_{Dup}(t,K)^2
-2\,\frac{\mathbb{E}^{d}\big[D^d_{0t}(\ell^d_t-\ell^f_t)\,\big|\,S_t=K\big]}{\mathbb{E}^{d}\big[D^d_{0t}\,\big|\,S_t=K\big]} .
\end{align*}
\end{prop}
\begin{proof}
By (\ref{eq: kernels}) and (\ref{eq: integratedSensitivities}), the two integrals in (\ref{eq: threeFactorGeneralMalliavin}) are the conditional expectations of the first components $Y^{d,S}_t=S_t\ell^d_t$ and $Y^{f,S}_t=S_t\ell^f_t$, which equal $K\ell^d_t$ and $K\ell^f_t$ on $\{S_t=K\}$. The foreign conditional expectation is taken under $\mathbb{Q}^{f^t}$; by Lemma \ref{lem: forwardMeasures} (ii) it equals the one under $\mathbb{Q}^{d^t}$. The second line is the usual change of numeraire between $\mathbb{Q}^{d^t}$, $\mathbb{Q}^{d^T}$ and $\mathbb{Q}^d$.
For completeness, the calibration equivalence can be established without
identifying the model density with the market density in advance. Write
$m(t,K)=\mathbb E^d[D^d_{0t}\delta(S_t-K)]=\partial_{KK}C_m(t,K)$ and
$c^i_t=r^i_t-r_i^0(t)$. Apply Lemma \ref{lem: duality} to smooth approximations of
$\mathbf1_{(K,\infty)}$. Dominated convergence handles the left sides;
the continuous weighted densities in (H3) handle the approximate identities
on the right sides. This gives
\begin{align*}
\mathbb E^d[D^d_{0t}c^d_t\mathbb I_{S_t>K}]
 &=K m(t,K)\mathbb E^{d^t}[\ell^d_t\mid S_t=K],\\
\mathbb E^d[D^d_{0t}S_tc^f_t\mathbb I_{S_t>K}]
 &=K^2m(t,K)\mathbb E^{d^t}[\ell^f_t\mid S_t=K].
\end{align*}
Here the first identity uses the domestic forward measure; the second uses
foreign duality followed by Lemma \ref{lem: forwardMeasures}. It\^o--Tanaka
therefore yields, for the model surface,
\[
\partial_t C_m=-(r_d^0-r_f^0)K\partial_K C_m-r_f^0C_m
+\tfrac12K^2\left(\sigma^2(t,K)+2\mathbb E^{d^t}
[\ell^d_t-\ell^f_t\mid S_t=K]\right)\partial_{KK}C_m.
\]
If $C_m=C$, division by the positive market density proves necessity.
Conversely, substituting (\ref{eq: mainFormula}) makes this precisely the
market forward Dupire PDE. The initial conditions coincide. At $K=0$,
$C_m(t,0)=\mathbb E^d[D^d_{0t}S_t]=S_0P^f_{0t}$ by the true-martingale
numeraire change, and $C_m(t,K)\to0$ as $K\to\infty$ by integrability of
$D^d_{0t}S_t$. The assumed uniqueness therefore implies $C_m=C$. This argument proves neither existence
nor uniqueness of the nonlinear calibrated SDE.

\end{proof}
Compared with (\ref{eq: singleMeasure}), the correction no longer contains $\partial^2_KC(t,K)$ nor tail expectations, and it is a conditional expectation of quantities carried by the particles themselves.
